% Book pp. 51-53 (PDF pp. 52-54)
\section{Sum of Sequences}

In year one we identified number patterns and classified them into arithmetic and
geometric sequences. While sequences are a list of numbers with defined patterns,
a series is the sum of the terms of a sequence. This year we will be finding the
sum of sequences and analysing the convergence or divergence of series.

If the general term of a series with $n$ term is known, then the complete series can
be written down in short notation as indicated by the following:
\begin{enumerate}
  \item $a + (a+d) + (a+2d) + \ldots + (a+[n-1]d) = \sum_{r=1}^{n} (a+[r-1]d)$
  \item $a + ar + ar^2 + \ldots ar^{n-1} = \sum_{k=1}^{n} ar^{k-1}$
  \item $1^2 + 2^2 + 3^2 + \ldots n^2 = \sum_{r=1}^{n} r^2$
  \item $-1^3 + 2^3 - 3^3 + 4^3 + \ldots (-1)^n n^3 = \sum_{r=1}^{n} (-1)^r r^3$
\end{enumerate}

\noindent\textbf{Notes:}
\begin{enumerate}
  \item[a.] It is sometimes more convenient to count the term of a series from zero
  rather than 1.

  For example:
  \[
    a + (a+d) + (a+2d) + \ldots + (a+[n-1]d) = \sum_{r=0}^{n-1} (a+rd) \text{ and}
  \]
  \[
    a + ar + ar^2 + \ldots ar^{n-1} = \sum_{k=0}^{n-1} ar^{k}
  \]
  In general, for a series with $n$ terms starting at $u_0$:
  \[
    u_0 + u_1 + u_2 + u_3 + \ldots u_{n-1} = \sum_{r=0}^{n-1} u_r
  \]
  \item[b.] We may also use the sigma notation for ``infinite series'' such as those already
  encountered in the sum to infinity of a geometric series.

  For example: $1 + \frac{1}{3} + \frac{1}{3^2} + \frac{1}{3^3} + \ldots = \sum_{r=1}^{\infty} \frac{1}{3^{r-1}}$
  or $\sum_{r=0}^{\infty} \frac{1}{3^{r}}$
\end{enumerate}

\begin{activity}{Activity 2.1 -- Sums of Sequences}
Perform this activity in pairs or groups.

In year one we learnt about Arithmetic and Geometric Progression. Let us
revise and build on that knowledge by creating the sequence terms and find
the sum of the terms of the expressions below.
\begin{enumerate}
  \item Evaluate $\sum_{n=0}^{6} 2(n-3)$
  \item Find the sum $\sum_{k=1}^{50} (3k+2)$
  \item Evaluate $\sum_{r=1}^{3} \left(\frac{1}{2}\right)^r$
\end{enumerate}

\solution
\begin{enumerate}
  \item $\sum_{n=0}^{6} 2(n-3) = 2(0-3) + 2(1-3) + 2(2-3) + 2(3-3) + 2(4-3) + 2(5-3) + 2(6-3)$
  \begin{align*}
    &= 2(-3) + 2(-2) + 2(-1) + 2(0) + 2(1) + 2(2) + 2(3) \\
    &= (-6) + (-4) + (-2) + 0 + 2 + 4 + 6 \\
    &= 0
  \end{align*}
  \item $\sum_{k=1}^{50} (3k+2)$
  \begin{align*}
    u_1 &= 3(1) + 2 = 5 \\
    u_{50} &= 3(50) + 2 = 152 \\
    s_k &= \frac{k\,(u_1 + u_{50})}{2} \\
    s_k &= \frac{50(5+152)}{2} = 3925
  \end{align*}
  \item $\sum_{r=1}^{3} \left(\frac{1}{2}\right)^r$
  \begin{align*}
    \text{when } r = 1,\ \left(\tfrac{1}{2}\right)^1 &= \tfrac{1}{2} \\
    \text{when } r = 2,\ \left(\tfrac{1}{2}\right)^2 &= \tfrac{1}{4} \\
    \text{when } r = 3,\ \left(\tfrac{1}{2}\right)^3 &= \tfrac{1}{8} \\
    \Longrightarrow \tfrac{1}{2} + \tfrac{1}{4} + \tfrac{1}{8} &= \tfrac{7}{8}
  \end{align*}
\end{enumerate}
\end{activity}

\subsection{Method of undetermined coefficients}

The method of undetermined coefficients is a powerful technique for finding the
general solution of linear recurrence relations with constant coefficients

If the $n$th term of a linear sequence is denoted by $x_n$, then the recurrence relation
is of the form: $x_{n+1} = f(x_n)$ for some of the function $f$.

For example, in $x_{n+1} = 4 - \frac{x_n}{2}$, when $x_2 = 2$, then $x_3$ is expressed as
$x_{2+1} = 4 - \frac{x_2}{2}$
\begin{align*}
  x_3 &= 4 - \frac{2}{2} \\
  x_3 &= 4 - 1 \\
  x_3 &= 3
\end{align*}

\begin{example}{Example 2.1}
The $n$th term of a sequence is $U_n$. If $U_1 = 3$ and $U_{n+1} = U_n + 5$, find $U_3$.

\solution
\begin{align*}
  U_{n+1} &= U_n + 5 \\
  U_1 &= 3 \\
  U_2 &= U_{1+1} = U_1 + 5, \text{ but } U_1 = 3 \\
  U_2 &= 3 + 5 = 8 \\
  U_3 &= U_{2+1} = U_2 + 5 \\
  U_3 &= 8 + 5 = 13
\end{align*}
\end{example}
